\chapter[East Asian Ritual and Popular Belief]{East Asian Rituals, the Three Teachings, and Popular Belief}
\section{A Stack of Spirit Money}
First, pick up an object: a stack of spirit money.

It might be coarse yellow paper stamped with blurred images of copper coins, or finely printed "Bank of the Underworld" notes with astonishing denominations starting at a trillion. A stack of several dozen sheets weighs almost nothing in your hand. It may cost less than one yuan, yet every spring and late summer, countless families across China, Taiwan, Hong Kong, Singapore, Vietnam, and Korea solemnly burn it, turning it into smoke, an image, a rising wisp of ash.

Watch this action for a moment: educated people, in the age of mobile payments, earnestly burning paper they know has no purchasing power. Ask why and you may receive several rather different answers: "For our relatives on the other side." "It's a gesture." "Doesn't everyone do it?" "Stop asking and get on with burning it."

This stack of paper is our doorway into the chapter. Behind it lies an enormous, ancient, still-functioning world: ancestral tablets, the earth god's shrine at the village entrance, Qingming and winter-solstice observances, large urban Buddhist temples, mountain Daoist monasteries, monks chanting at funerals, and auspicious dates chosen for marriage documents. Sometimes these are called "superstition," sometimes "traditional culture," sometimes "the Three Teachings: Confucianism, Buddhism, and Daoism." Which name gets them right? Or which part does each get right?

\section{A Family on the Hillside on Qingming Morning}
Come with me to a particular scene.

Time: nine in the morning on April 5, Qingming, the Tomb-Sweeping Festival. Place: a low hill outside a county town in southern China. Gravestones crowd its gentle slope like a quiet little housing estate. Rain fell last night. The earth is dark brown-black and soft underfoot; mud soon coats the sides of your shoes. Three smells mingle in the air: wet grass, the smoke left by firecrackers, and the scorched scent of paper burning elsewhere.

A family approaches along the hillside. Grandmother leads, in her seventies, carrying a bamboo basket containing a cooked chicken, a plate of steamed sponge cake, three apples, a bottle of rice wine, and that stack of spirit money. Behind her comes her son, the father of fifteen-year-old A-Yuan, carrying a hoe to clear weeds in front of the grave. An aunt carries paper models of "household necessities": a phone and a house, quite carefully made. A-Yuan trails behind, reluctantly holding chrysanthemums that his mother insisted on buying: "Civilized grave-visiting---that's what your teacher says."

At the grave, a whole sequence begins. Father hoes away the weeds. Aunt wipes the photograph on the gravestone with a cloth she brought: Grandfather in a Zhongshan suit, looking solemn. Grandmother arranges the chicken with its head toward the stone, pours three cups of wine and spills them on the ground, then lights three sticks of incense and distributes incense to everyone, asking them to "speak to Grandpa in your heart." Then comes the paper burning. Her hand is steady as she strikes a match and feeds the notes into the fire one by one. She murmurs something indistinct; only a few words can be heard: "... protect A-Yuan ... his studies ... keep him safe ..."

Smoke rises and drifts sideways. A-Yuan coughs and steps back half a pace.

At that moment he cannot hold back the question he has wanted to ask for so long: "Grandma, do you really believe Grandpa can receive this money?"

The hillside falls quiet for a second. Father stops hoeing. Aunt glances at Grandmother. Without looking up, Grandmother adds another sheet to the fire. After a while she says, "Whether I believe or not, burning it sets my mind at ease."

Stand on this hillside for a moment too. Nobody in this ceremony seems ignorant, yet nobody can explain it completely. This family will stay with us throughout the chapter.

\section{What Are They Really Seeking?}
Let us first set out the tension, without rushing to an answer.

A-Yuan's question is really a chain of questions. First: Grandfather has died; can the dead still "receive" things? Science class gives a straightforward answer: no. Consciousness ceases after death, and burned paper becomes ash; no channel deposits it into an account. By that standard, the family's actions rest on a mistaken belief. Identifying a false belief is precisely the usual job of the word "superstition."

Second: Grandmother's answer is odd. She did not say he could receive it, nor that he could not. She said burning it set her mind at ease. Her concern seems to be something other than the paper's delivery arrangements. So for whose benefit is it burned? Grandfather's? Her own, as someone still living? The whole family's, gathered there? Or an inexpressible feeling that something would be wrong if it were not burned?

The third layer goes deeper. If anything unscientific should stop, why did the school's most scientifically minded physics teacher also burn incense at a temple before his child's college entrance examination? Why does A-Yuan's mother insist on "civilized grave-visiting" yet spend ages consulting the traditional almanac for a wedding date, refusing any date containing four? A gap separates what this family---and this whole society---says from what it does.

The real question emerges: \textbf{What exactly is someone doing when they light a stick of incense?} Stating a belief about the afterlife? Making a transaction with a deity? Expressing an unspoken longing? Fulfilling a family member's duty? Or doing several things at once, without clearly distinguishing them?

A larger question follows. Scenes like this have recurred across East Asia for more than two thousand years. States established special offices and grand ceremonies for them; scholars wrote mountains of arguments about them; Buddhism and Daoism competed over them and borrowed from one another. If we dismiss all this with the word "superstition," what are we dismissing: other people's ignorance, or our own ability to understand human beings?

Before reading on, take a side on the hillside. Does A-Yuan's question make sense to you, or does Grandmother's action? Remember your answer; we will revisit it at the end.

\section[Emperors, Priests, Monks, and Gods]{Emperors, Daoist Priests, Buddhist Monks, and the Earth God}
The family's incense does not stand alone. It is planted in an enormous web stretching from the emperor's altar to the village earth god's shrine. To hear this web's many voices, we must pull the camera back.

\textbf{The first voice: the state.} Beijing still has the Temples of Heaven, Earth, Sun, and Moon, along with the Imperial Ancestral Temple. Under imperial rule, sacrifice to Heaven was the emperor's exclusive privilege. An ordinary person performing it committed a grave offense, because the right to communicate with Heaven was itself part of imperial authority. The state also maintained something like a "household register" of deities. Local officials were required to offer timely sacrifices to gods on the recognized list, the "authorized cults." Gods outside it belonged to "unauthorized cults," called yinsi, where yin means excessive or transgressive; officials could demolish their temples. Confucius was incorporated too: successive emperors conferred additional titles on him, and sacrifices to him were major state ceremonies. Notice the implication: in East Asian tradition, \textbf{sacrifice was never merely a private spiritual matter; it was also part of political order}. Who could worship whom, and with what ceremonial rank, were questions of power.

\textbf{The second voice: the scholar.} Confucian scholar-officials had a subtle attitude toward spirits and gods. The Analects records a student asking what wisdom means. Confucius's answer, broadly, was to attend to human duties and respect spirits while keeping one's distance. Another student asked about death. Confucius replied, in effect: before understanding life, why rush toward the world after death? This is not atheism. Confucius never denies spirits' existence; he insists on turning attention back to human ethics: honor parents, conduct funerals carefully, and perform sacrifices on time. Their meaning concerns how the living conduct themselves, not whether spirits sign for delivery. For the scholar-official, sacrifice is a practice in becoming the kind of person one ought to be.

\textbf{The third voice: the Daoist priest.} Daoism offers another picture of the world: a system of gods and immortals between Heaven and Earth, including Wenchang, who oversees examinations; the City God, responsible for a city; and Mazu, responsible for a stretch of sea. People engage with this system through talismans, ritual registers, purification rites, offerings, and cultivation. When someone falls ill, a Daoist priest can be invited to expel harmful influences; before a ship sails, people visit Mazu's temple for protection. This is a set of professional techniques: Daoist priests are licensed intermediaries between humans and the divine world.

\textbf{The fourth voice: the Buddhist monk.} Buddhism arrived from India with things previously absent from East Asia: rebirth, karmic consequences, and a complete set of rites to assist the dead. Where do people go after death? Buddhism offered extremely detailed answers and accompanying solutions: inviting monks to chant and generate merit could help departed relatives fare better in rebirth. One source of the Ghost Festival, the Ullambana assembly, derives from the Buddhist story of Mulian rescuing his mother. With the combined power of monks, a son saves his mother from an unfortunate realm of existence. This theme---\textbf{rescuing one's own mother}---immediately joined with Chinese filial duty, securing Buddhism a permanent seat at the table of ancestral observance that mattered so much to Chinese people.

\textbf{The fifth voice, with the most speakers: ordinary people.} Their position is almost disarmingly practical: worship whichever deity proves effective. For a child's examination, honor both Wenchang and Confucius, and offer candles before the Buddha as well. When someone falls ill, see a doctor first, but also ask a Daoist priest for a talisman while an elderly relative murmurs to the Kitchen God. At a funeral, Buddhist monks chant to aid the deceased, Daoist priests perform rites to open the way, a diviner selects the burial hour, and the family wears traditional mourning clothes according to Confucian rites. \textbf{All Three Teachings attend the same funeral, and nobody feels an apology is owed.}

Another schedule weaves these voices together: \textbf{seasonal festivals}. At New Year, honor ancestors and welcome gods; at the Lantern Festival, light lamps; at Qingming, sweep graves; at the Dragon Boat Festival, ward off disease with mugwort and scented sachets, with commemorations of Qu Yuan added in later layers; at the Ghost Festival, assist wandering souls; at Mid-Autumn, honor the moon and reunite; at Double Ninth, climb heights and respect elders; at the winter solstice, honor ancestors and welcome winter. Seven or eight major points divide the year, each prescribing whom to honor, what to eat, and where the family should gather. Festivals form popular religion's calendrical skeleton. Temple deities have birthdays, family ancestors have death anniversaries, and state rites have appointed dates. Three schedules ultimately combine in one almanac. A person in traditional society need not "believe" anything: simply living through the year involves a whole religious life.

This is what most puzzles outsiders about East Asian religion. In European and West Asian traditions, religious identity resembles nationality: a Christian is not a Muslim, conversion is momentous, and exclusivity is written into doctrine. Yet a traditional Chinese person---a Qing merchant, say---might instruct his son before Confucius's image in the morning, invite monks to chant for his deceased mother at midday, and consult a Daoist about his shop's feng shui in the evening. He neither thinks he believes in three religions nor can clearly say which he belongs to. Confucianism, Buddhism, and Daoism resemble not three sharply bounded modern religious organizations but three coexisting services or complementary languages: Confucianism handles human order and ethics, Buddhism life, death, and the afterlife, and Daoism cosmic operations and particular affairs. They competed for imperial and popular "market share," with periods of Buddhist suppression and Daoist favor, but more often coexisted in the same city, family, or even temple. In many Chinese temples, images of Confucius, Laozi, and Shakyamuni sit side by side, all receiving plentiful incense.

This coexistence is not uniquely Chinese. In Japan, Buddhist temples and Shinto shrines coexisted for more than a millennium. People registered births at shrines, married with Shinto rites, and entrusted funerals to Buddhist temples; Japanese does not even have a word corresponding to Western "conversion." On the Korean peninsula, Confucian ceremonies and Buddhist temple bells coexisted for centuries. Even today, a Korean family may perform Confucian ancestral rites, charye, while also containing Christians. Whether to bow to ancestors remains a real debate in Korean churches. Vietnamese temples offer incense to the Buddha, Confucius, the Jade Emperor, and local mother goddesses together. Look farther away to ancient Rome: households honored their own domestic and hearth deities, the state maintained official ceremonies, and newly conquered regions' gods were invited into the Roman pantheon. More was better; none demanded exclusivity. \textbf{Imagining religion as an either-or membership card belongs to only a few of the world's traditions; imagining it as a toolbox with different tools for different tasks comes closer to the norm across most places and periods of human history.}

\section[Thinking Bridge: Practice Before Belief]{Thinking Bridge: Ask What People Do Before Asking What They Believe}
For that incense on the hillside, we need a portable tool. It is simple enough to state in one sentence: \textbf{split "religion" into three independent questions and ask them separately}.

\begin{itemize}
\item Question one: what does this person \textbf{believe}? (Doctrine and belief.)
\item Question two: what does this person \textbf{do}? (Ritual and practice.)
\item Question three: to whom does this person \textbf{belong}? (Organization and identity.)
\end{itemize}
Burning spirit money puzzles us because we assume these questions must come bundled: first belief, then action, then identity. "Because I believe ancestors can receive it, I burn paper, and therefore I am a follower of such-and-such." This intuition grows from experiences of religions such as Christianity and Islam, where creeds stand at the beginning of scripture, entry has clear rituals, and membership is explicit. It is powerful, but it is a locally made pair of glasses, not a universally reliable pair of eyes.

Apply the three questions to the hillside. What does Grandmother \textbf{believe}? A standard answer is hard to obtain; she might say, "Who knows? Better to believe it exists." What does she \textbf{do}? That is clear: arrange offerings, light incense, burn paper, bow, year after year. To which organization does she \textbf{belong}? None; she is on no membership register. The three questions receive three different answers, yet the ceremony has continued for two thousand years. This suggests that \textbf{in East Asian religious life, doing is central, belief is an optional background, and belonging is often absent altogether}. Scholars call creed-centered religion "orthodoxic" and traditions centered on correct practice "orthopraxic." You need not memorize these terms; remember the three questions.

The tool can be expanded into a five-question ritual checklist: \textbf{Who participates? When and where? Whom do they address? What objects and actions do they use? What changes afterward?} Apply it to Qingming grave-visiting. Participants are a family related by blood; absentees attract comment. The time is seasonal, Qingming or the winter solstice, determined by the calendar rather than mood. The recipients are particular ancestors in family graves, not abstract gods. Objects include food, wine, incense, and paper, all ritualized versions of everyday things. What changes afterward is not the dead person's situation---even Grandmother may not insist on that---but the living people's condition: the family has assembled, seniority has been rehearsed, and some debt of gratitude has been discharged. Without asking about belief, simply observing actions reveals the mechanism's gears one by one.

Next time you encounter a seemingly inexplicable ritual---moving a graduation tassel, an Olympic champion biting a gold medal, a gaming team's fixed pre-match gesture---do not rush to judge. First apply the three questions and then the five. Your eventual judgment will have a different flavor.

\section{Confucius on Sacrificing "As If Present"}
Now let us read a short primary source. One of the most discussed statements about sacrifice from more than two thousand years ago appears in the Analects, "Eight Rows of Dancers":

\begin{quote}
He sacrificed as if the ancestors were present; he sacrificed to spirits as if the spirits were present. The Master said, "If I do not take part in the sacrifice, it is as though I have not sacrificed."
\end{quote}
Literally, this means treating ancestors as truly present when sacrificing to them, and spirits as truly present when sacrificing to spirits. Confucius says that if he cannot participate personally, it is as though no sacrifice took place.

Weigh the words carefully; their subtlety is remarkable. "As if the spirits were present": Confucius does \textbf{not} say they are present, but \textbf{as if} they are. That small phrase quietly shifts the center of gravity from the divine side to the human side. What matters is not whether spirits actually arrive, a question on which he declines to take a position, but whether the participant inwardly treats them as present. The next sentence strengthens the point: without his own participation, it is as though there were no sacrifice. If sacrifice merely delivered goods to spirits, hiring someone to arrange delivery should work equally well. Yet Confucius says that personal absence means no sacrifice. The real recipient, then, is the participant: you arrive, perform the rites, and deliberately exercise respect. That is what counts.

In the Analects, "Learning," Confucius's student Zengzi offers another statement:

\begin{quote}
Attend carefully to the end of life and remember those long departed; the people's virtue will grow deep and generous.
\end{quote}
Its general meaning is that carefully conducting parents' funerals and reverently remembering distant ancestors will make people's character more sincere and generous. Notice the chain of reasoning: sacrifice is the means; public virtue is the end. Zengzi is concerned not with how many offerings ancestors receive but with \textbf{how repeated practice in remembering one's roots trains gratitude and repayment of kindness into a society's muscle memory}. Someone who conscientiously remembers deceased parents will probably not treat the living too harshly. This is Confucianism's central defense of sacrifice, grounded throughout in human life.

The Analects, "Governing," adds a related rule. Discussing filial duty, Confucius says, broadly, that parents should be served according to ritual while alive, buried according to ritual after death, and thereafter honored according to ritual. One ritual order connects life, death, and sacrifice. Taken together, the Confucian position is clear: \textbf{ancestral sacrifice extends family ethics; it is not a receiving station for supernatural events}. This is less distant than it first appears from Grandmother's "Burning it sets my mind at ease."

\section{One Stick of Incense, Several Explanations}
We can now offer several genuinely different explanations of the hillside scene. First distinguish four things: what happened, a family visiting a grave and burning paper, which is uncontroversial at the evidential level; why this ritual has lasted two thousand years, where explanations clash; how participants understand it, as when Grandmother says it gives peace of mind; and how we evaluate whether it should continue, a value judgment. What follows compares the second level.

\textbf{Explanation one: it addresses the psychology of the living.} Grief needs an outlet; longing needs an address. After someone dies, the survivors' impulse to keep doing something for them has no recipient. Ritual supplies one. Psychologists would point out that fixed rituals, performed on the same day with the same actions each year, can provide certainty against death's experience of lost control. This explanation directly takes up Grandmother's peace of mind without needing to settle whether spirits exist, making it scientifically safe. Its limitation is equally clear: why does the ritual take precisely this form? Why food and money, why Qingming rather than a random day, and why must the whole family attend? If peace of mind were all that mattered, would lighting a candle alone at home not be easier?

\textbf{Explanation two: it assembles the family and community.} Anthropologists supply this missing piece: the ritual's true product is the group itself. Qingming summons relatives scattered across cities to one hillside. The entire family sees who attends, who is absent, and who bows. It is an annual membership roll call and inspection of generational order. At the local level, City God temples, Mazu temples, and earth god shrines form a \textbf{temple network}. A temple is not merely a place to burn incense but a meeting hall, charity, festival committee, and dispute-mediation office. Historically, coastal Mazu temple associations relieved disaster victims, mediated merchant disputes, and accommodated travelers. A village's rotation of temple-festival duties rehearses cooperation in public affairs generally. This persuasive explanation reveals the social organization behind ritual. Its limitation is that it can make participants seem like pieces manipulated by functions, as though Grandmother went grave-visiting to integrate the family. She herself is thinking about Grandfather. Replacing participants' motives with an observer's inferred functions can distort the picture.

\textbf{Explanation three: it is an exchange between humans and gods.} This account treats ritual as a transaction: people offer gifts, incense, and respect; deities return protection, effective help, and reassurance. A responsive deity attracts more incense; an unresponsive one sends worshippers to another temple. Ordinary people's apparent inconsistency becomes astute consumer logic. This captures real bargaining in popular religion---making and fulfilling vows resembles contractual language---and explains competition among temples. But reducing everything to advantage has limits. Why does A-Yuan's mother consult the almanac despite knowing that a wedding date does not determine a marriage's quality? No concrete benefit comes from that transaction. She seeks relief from the feeling that omitting the action would be wrong. A pure exchange model cannot account for that kind of unease.

\textbf{Explanation four: it is an ethics lesson without a classroom.} This is Confucianism's own account, already voiced by Zengzi: careful funerals and remembrance teach the living gratitude, responsibility, and generational respect. A child learns filial duty not only from a textbook but at the moment someone hands them incense on the hillside. This explanation centers participants' own understanding, using reasons they articulated two thousand years ago. Its weakness is excessive respectability: it can overlook less dignified ingredients such as fear, calculation, and competitive display. Nor does it explain why Buddhists and Daoists, with completely different worldviews, can peacefully chant their own scriptures at the same funeral.

Each explanation holds part of the truth. Here we must especially watch for \textbf{a counterexample that is easily misread}. Many people, including East Asian intellectuals impatient for modernization, glanced at temple networks and concluded that demolishing these centers of superstition could only benefit society. Yet investigating a Mazu or City God temple's accounts and personnel often revealed substantial public functions: charity, mediation, deliberation, and relief. Demolition removed not only a divine image but a community's living room. The object was misidentified because observers asked only what people believed, not what they did. Twentieth-century societies paid a real price for that mistake.

\section{Further Challenge: Religion Without Walls}
We can now introduce the chapter's weightiest theory, from the Chinese sociologist C. K. Yang. In the mid-twentieth century, his classic Religion in Chinese Society challenged Western scholars with an unsettling question: you keep counting Chinese believers who resemble Christians, find small numbers, and conclude that Chinese people have no religion. Might you have asked the wrong question from the outset?

Yang's approach distinguishes two forms. \textbf{Institutional religion} has independent scriptures, clergy, organizations, and membership lists, like a courtyard with walls and a gate. Chinese Buddhist and Daoist monasteries belong here and are indeed limited in scale. \textbf{Diffused religion} has no walls, gate, or register. Its doctrines dissolve into ethical textbooks; household and lineage heads double as clergy; its churches are ancestral halls, graves, and the space before the kitchen stove; its rituals mingle with weddings, funerals, and seasonal observances. Less like a building than air, it permeates family, lineage, merchant guild, and state ceremony. Followers of the first can be counted; participants in the second approach the entire population.

This distinction immediately transforms old questions. Do Chinese people believe in religion? By institutional measures, most belong to nothing; by diffused measures, almost everyone participates---visiting graves at Qingming, sharing lucky-carp images online, or avoiding unlucky words before an examination. The century-old knot of whether Confucianism is a religion also loosens. As a system with classics and sacrificial rites, deeply embedded in state ritual and family ethics, Confucianism does not fit a Western mold centered on doctrine and organization. Through diffused religion, however, it appears as East Asia's most thoroughly pervasive example. \textbf{The opposing sides are using different rulers, borrowed from someone else's factory.}

This theory's boundaries also need testing. If religion becomes as broad as air, does everything become religion? New Year? A school opening ceremony? If all count, does the word lose its boundaries? Scholars genuinely debate these questions without a common answer. Even without settling them, the theory teaches at least this: \textbf{when a question repeatedly fails to make sense of something, inspect not only the answer but the default mold hidden inside the question}. Asking only whether something is superstition cannot explain the hillside family, just as counting registered believers cannot explain a civilization that makes ethics, ancestors, and seasonal observances part of everyday life.

\section[Digital Rituals and Exam Incense]{Electronic Wooden Fish and Incense Before Examinations}
This two-thousand-year-old question is on your phone right now.

Open an app store and download an "electronic wooden fish": tap the screen and gain one unit of merit. Originally a self-mocking joke among young people, it becomes more serious with use; on the day before postgraduate entrance examinations, tapping surges. A lucky-carp image circulates millions of times on social media with the caption "Share this and you'll pass." During college entrance examination week, Confucian temples and Wenchang pavilions enter their annual busy season, with parents queuing to touch statues symbolizing examination success. That same week a class teacher says, "Superstition is useless; check your answers carefully," then quietly drinks "lucky tea" in the staff room. At Qingming, people unable to return home light virtual candles and offer virtual flowers on online memorial platforms. Spirit money has not disappeared; it has moved onto servers. In many cities, temples become crowded weekend destinations for young people. Media call it an "incense-burning craze," while participants speak plainly: "It's not belief. I have nowhere else to put my anxiety."

How should we interpret these phenomena? First, a serious exercise: \textbf{state the strongest possible case for the position "This is all superstition and should be criticized."} This book repeatedly asks us to explain opponents' reasons so fully that they themselves would agree, then decide whether we accept them.

That case has at least three layers. First, cognition: no causal mechanism links sharing a carp image to examination marks. Spending time on it, even at psychological cost alone, invests in the false worldview that wishes run the world. That worldview may collect interest in more consequential settings, such as trusting folk remedies over doctors when ill. Second, economics: the incense economy contains real exploitation---outrageously priced incense, fortune-telling businesses, and false masters profiting from anxiety. Superstition lubricates this harvesting machine. Third, responsibility: attributing results to lucky carp shifts blame for inadequate study onto fortune, corroding responsibility for one's own life. None of these arguments is weak; honesty requires facing them directly.

But the other side deserves a full hearing too. Few electronic wooden-fish users truly think their phones connect to a cosmic merit system. They do something older: when facing something important beyond their complete control---an examination, job search, or relative's illness---they perform a low-cost action to loosen anxiety's grip. It is Grandmother's "Burning it sets my mind at ease," decades later. A virtual memorial candle is not a parcel sent to the dead but an address for the living person's longing, only a network cable away from Confucius's "as if present." Opposing prayer to responsibility may also be unfair: children who touch lucky statues usually still check their examination papers, and incense-burning jobseekers still submit applications. \textbf{Human beings have never been an either-or choice between rationality and superstition; they work hard where they have control and light incense where they do not.} A soldier honoring the war god two thousand years ago could still sharpen his blade during the day.

A different danger deserves attention. When we mock elderly people burning paper while sharing lucky carp ourselves, the line labeled "superstition" often separates not truth from falsehood but us from them. On our side it is emotional regulation; on theirs, backward ignorance. Yet the three-question tool reveals almost identical ritual structures. Perhaps this is the tool's greatest contemporary use: helping you see what you yourself are doing before deciding how to judge others.

\section[Without Belief, Would You Still Bow?]{Your Question: Without Belief, Would You Still Bow?}
Back on the hillside, the paper has burned away, leaving red ash. The family begins collecting the offerings. Custom says the chicken and sponge cake should be taken home and shared: the living eat food the ancestors have "enjoyed," and the meal itself belongs to the ritual. Chewing a chicken leg, A-Yuan remembers his question.

Now it is yours, in a form harder to evade:

Suppose you are completely certain---not merely saying so, but genuinely and thoroughly certain---that dead people receive nothing and gods do not exist. Next Qingming, would you still climb the hill? Would you still set out the chicken and light those three sticks of incense?

Before answering, count your cards again:

If you say, "No, because it does nothing," respond to Zengzi: perhaps ritual was never delivery for the dead but an ethical lesson and psychological anchor for the living. What you cancel may be not a superstitious event but the year's only full family gathering, its only serious discussion of where you come from. How would you replace that loss? Or do you think no replacement is needed?

If you say, "Yes, for peace of mind," respond to the critics: is knowingly performing an ineffective action a kind of dishonesty with yourself? Does "everyone does it" justify doing something? When someone sends a letter knowing there is no recipient, is that deep feeling or self-deception? Where is the boundary?

There is a further twist. Is Confucius's "as if present" still defensible today, asking not whether you believe but only that you treat the recipient as present while acting? Is it a third path invented two thousand years ago, or an elegant way of hiding the question?

There is no standard answer. But note one fact: however you answer, you speak within an unbroken discussion more than two thousand years long. Confucius and Zengzi discussed it; persecuting emperors and protective monks argued over it; radical young temple-demolishers and elderly temple-keepers clashed over it; those sharing lucky carp and those mocking them continue it. The smoke on the hillside has not dispersed. Only the generations looking up at it keep changing.
